% !TeX root = ../../Book2.tex
% 纲要标题：### 16.2 双模实现的函子
% 纲要位置：计划纲要.md，第 1570 行。

\section{双模实现的函子}
\label{b2:sec:1602}

为了把两个模范畴联系起来，需要一个同时接受两侧标量作用的模。先从任意 \((R,S)\) 双模 \(P\) 出发，写出张量函子、Hom 函子及它们的伴随映射；然后再把 \(S\) 选为 \(P\) 的自同态环的反环。这样可以分别检查构造本身的合法性，以及何种条件使它成为等价。

% 纲要要点（供目录核对）：
%
% * 张量函子与 Hom 函子；调用第7.1节的有限投射 Hom–张量同构
% * 自同态环和反环约定；逐式核对预复合和右作用，正式证明任意环上的行列互逆双模及一般模上的余单位
% * 伴随的单位及余单位
%

\subsection{张量函子与 Hom 函子}
\label{b2:subsec:1602:01}

设 \(R,S\) 为含幺环，给定幺性双模 \({}_RP_S\)，定义
\[
F(M)=\operatorname{Hom}_R(P,M),\qquad
G(N)=P\otimes_SN.
\]
这里 \(M\) 为幺左 \(R\) 模，\(N\) 为幺左 \(S\) 模，因此 \(F:R\text{-}\mathbf{Mod}\to S\text{-}\mathbf{Mod}\)、\(G:S\text{-}\mathbf{Mod}\to R\text{-}\mathbf{Mod}\)。\(G(N)\) 的左 \(R\) 作用留在第一个因子；\(F(M)\) 的左 \(S\) 作用由
\begin{equation}\label{b2:eq:morita-hom-action}
(sf)(p)=f(ps)
\end{equation}
给出：\(s\) 先在输入端作用，再由 \(f\) 求值，因而这是预复合。它仍为左 \(R\) 线性映射，因为 \(P\) 的两侧作用相容；而
\[
\bigl(s_1(s_2f)\bigr)(p)=f\bigl((ps_1)s_2\bigr)
=\bigl((s_1s_2)f\bigr)(p),
\]
所以确为左 \(S\) 作用。

\ProseParagraphBreak
对 \(a:M\rightarrow M'\)，令 \(F(a)(f)=af\)；对 \(b:N\rightarrow N'\)，令 \(G(b)=\operatorname{id}_P\otimes b\)。复合与恒等关系由同态复合和张量函子性得到。它们都是加性函子。由\cref{b2:thm:hom-tensor-adjunction}，有关于两变量自然的对应
\[
\operatorname{Hom}_R(P\otimes_SN,M)
\cong\operatorname{Hom}_S\bigl(N,\operatorname{Hom}_R(P,M)\bigr),
\]
把 \(h\) 送到 \(n\mapsto\bigl(p\mapsto h(p\otimes n)\bigr)\)。因此 \(G\) 左伴随于 \(F\)。一般地，\(F\) 左正合，\(G\) 右正合；\(P\) 左投射时，\cref{b2:cor:projective-hom-exact}才使 \(F\) 正合。

\ProseParagraphBreak
若这里的双模 \({}_RP_S\) 作为左 \(R\) 模有限生成并投射，\cref{b2:prop:finite-projective-hom-tensor}便把已经定义的 \(F(M)=\operatorname{Hom}_R(P,M)\) 写成张量函子。具体地，令 \(P^*=\operatorname{Hom}_R(P,R)\)，它接受 \((S,R)\) 双模作用
\[
(s\lambda)(p)=\lambda(ps),\qquad
(\lambda r)(p)=\lambda(p)r.
\]
该命题给出关于每个左 \(R\) 模 \(M\) 自然的左 \(S\) 模同构
\[
P^*\otimes_RM\longrightarrow F(M),\qquad
\lambda\otimes m\longmapsto\bigl(p\mapsto\lambda(p)m\bigr).
\]
目标的左 \(S\) 作用正是式\eqref{b2:eq:morita-hom-action}中的作用，因此这个同构保留当前的双模侧别。公式中的标量 \(\lambda(p)\) 作用在 \(m\) 左侧；交换这两个符号在一般环上并没有意义。

\subsection{自同态环和反环约定}
\label{b2:subsec:1602:02}

现在固定一个含幺环 \(R\) 及幺左 \(R\) 模 \(P\)，取
\[
\Delta=\operatorname{End}_R(P),\qquad S=\Delta^{\mathrm{op}}.
\]
\symidx{\(S=\operatorname{End}_R(P)^{\mathrm{op}}\)：投射生成元对应的自同态反环}
若 \(s\in S\) 对应自同态 \(\alpha_s\in\Delta\)，定义 \(ps=\alpha_s(p)\)。因为
\[
\alpha_{st}=\alpha_t\circ\alpha_s,
\qquad (ps)t=\alpha_t\bigl(\alpha_s(p)\bigr)=p(st),
\]
这是右 \(S\) 作用；各 \(\alpha_s\) 左 \(R\) 线性，保证它与左作用相容。因此 \(P\) 成为 \((R,S)\) 双模，上一个小节的函子与伴随立即适用。

两侧作用可记为
\[
{}_RP_S,\qquad {}_SP^*_R,\qquad
P^*=\operatorname{Hom}_R(P,R),\qquad
S=\operatorname{End}_R(P)^{\mathrm{op}}.
\]
自同态的复合在右 \(S\) 作用中反序一次，因而这里必须取反环。

\ProseParagraphBreak
在此约定下，\(F(P)=\operatorname{End}_R(P)\) 作为左 \(S\) 模自然同构于正则左模 \(S\)：将 \(\alpha_s\) 对应于 \(s\)。事实上，\(t\alpha_s\) 的作用是先做 \(\alpha_t\) 再做 \(\alpha_s\)，即 \(\alpha_s\alpha_t=\alpha_{ts}\)，正对应左乘 \(t\)。所以在将来由这些函子给出的等价中，与正则左 \(S\) 模对应的是选定的 \(P\)，未必是正则左 \(R\) 模。这是一个左模识别；作为环，\(\Delta=S^{\mathrm{op}}\)，不能因此省去反环。

\ProseParagraphBreak
最基本的例子是 \(P={}_RR\)。每个左线性自同态都是右乘某个 \(a\in R\)，所以 \(\Delta\cong R^{\mathrm{op}}\)、\(S\cong R\)，Hom 在一处求值后恢复原模。若取左自由模 \(P=R^{1\times n}\)、\(n\ge1\)，写成行向量方便记录右作用；非交换系数不能随意搬到另一侧，所以不能把以下公式直接改成左乘列：每个左线性自同态唯一为右乘矩阵 \(B\in M_n(R)\)，两次右乘的算子复合满足
\[
D_BD_C(p)=(pC)B=p(CB).
\]
所以 \(S\cong M_n(R)\)。映射 \(f:P\rightarrow M\) 对应其在各标准行上的值组成的列
\[
\bigl(f(e_1),\ldots,f(e_n)\bigr)^{\mathsf T}\in M^n.
\]
由 \(e_iB=\sum_j b_{ij}e_j\)，左 \(S\) 作用恰为矩阵乘这个列。这便与\cref{b2:prop:matrix-morita-equivalence}的显式等价相合。

\subsection{伴随的单位及余单位}
\label{b2:subsec:1602:03}

对任意 \((R,S)\) 双模，伴随的单位和余单位由\cref{b2:thm:hom-tensor-adjunction}的两种求值公式给出：
\begin{align}
\eta_N:N&\longrightarrow FG(N),
&\eta_N(n)(p)&=p\otimes n,\label{b2:eq:morita-unit}\\
\varepsilon_M:GF(M)&\longrightarrow M,
&\varepsilon_M(p\otimes f)&=f(p).\label{b2:eq:morita-counit}
\end{align}
第一式比较先取张量再取 Hom 所得对象与原来的 \(N\)，第二式则将一个输入 \(p\) 与同态 \(f\) 配对，尝试恢复原来的 \(M\)。第二式的平衡性为 \(f(ps)=(sf)(p)\)，第一式的左 \(S\) 线性为 \(p\otimes sn=ps\otimes n\)。相应的左 \(R\) 线性分别来自 \(P\) 的左作用和 \(f\) 的线性。因此两者都是所在模范畴中的同态，不只是集合映射。

\begin{proposition}\label{b2:prop:morita-adjunction-maps}
设 \(R,S\) 为含幺环，\({}_RP_S\) 为幺性 \((R,S)\) 双模，令 \(F=\operatorname{Hom}_R(P,-)\)、\(G=P\otimes_S-\)。对左 \(S\) 模 \(N\) 和左 \(R\) 模 \(M\)，定义
\[
\eta_N:N\longrightarrow FG(N),\quad n\longmapsto\bigl(p\longmapsto p\otimes n\bigr),
\qquad
\varepsilon_M:GF(M)\longrightarrow M,\quad p\otimes f\longmapsto f(p).
\]
这些映射关于 \(N,M\) 自然，并满足
\begin{equation}\label{b2:eq:morita-triangles}
\varepsilon_{G(N)}G(\eta_N)=\operatorname{id}_{G(N)},\qquad
F(\varepsilon_M)\eta_{F(M)}=\operatorname{id}_{F(M)}.
\end{equation}
若 \(S=\operatorname{End}_R(P)^{\mathrm{op}}\)，且 \(P\) 的右 \(S\) 作用取自其左 \(R\) 线性自同态，则 \(\eta_S\) 与 \(\varepsilon_P\) 都是同构，不要求 \(P\) 已是生成元。
\end{proposition}
\begin{proof}
若 \(b:N\rightarrow N'\)，两种从 \(N\) 到 \(FG(N')\) 的路径在 \(n,p\) 处都为 \(p\otimes b(n)\)；若 \(a:M\rightarrow M'\)，余单位的两种路径在 \(p\otimes f\) 处都为 \(a\bigl(f(p)\bigr)\)，所以自然。第一条三角恒等式把 \(p\otimes n\) 送到 \(\eta_N(n)(p)=p\otimes n\)；第二条把 \(f\) 送到函数 \(p\mapsto\varepsilon_M(p\otimes f)=f(p)\)，故都成立。

在所给自同态环约定下，识别 \(P\otimes_SS\cong P\) 及 \(F(P)\cong S\)。\(\eta_S(s)\) 对应 \(p\mapsto ps=\alpha_s(p)\)，再对应到 \(s\)，所以是恒等映射。\(\varepsilon_P\) 在这些识别下为 \(P\otimes_SS\rightarrow P\)、\(p\otimes s\mapsto ps\)，是张量积的单位同构。
\end{proof}

这些特殊分量为同构，并不说明其他模上的分量也是同构。为了在全部模上得到同构，需要把模写成适当直和之间映射的余核，并证明函子保留相应的直和与正合尾段。有限生成、投射和生成元三个条件将在下一节分别承担这些任务。

\begin{proposition}[行列双模的互逆配对]\label{b2:prop:matrix-row-column-bimodules}
设 \(R\) 为含幺环、\(n\ge1\)，令 \(S=M_n(R)\)、\({}_RP_S=R^{1\times n}\)、\({}_SQ_R=R^{n\times1}\)，各作用由矩阵乘法给出。令 \(u_i\)、\(v_i\) 分别为第 \(i\) 个标准行、标准列，则矩阵乘法诱导双模同构
\[
P\otimes_SQ\longrightarrow R,\quad p\otimes q\longmapsto pq,
\qquad Q\otimes_RP\longrightarrow S,\quad q\otimes p\longmapsto qp.
\]
它们的逆分别为 \(r\mapsto ru_1\otimes v_1\) 和
\[
(b_{ij})\longmapsto\sum_{i,j}v_i b_{ij}\otimes u_j.
\]
对任意幺左 \(R\) 模 \(M\)，将 \(\operatorname{Hom}_R(P,M)\) 识别为左 \(S\) 模 \(M^n\) 后，余单位也是同构：
\[
P\otimes_SM^n\longrightarrow M,\quad
(p_i)_i\otimes(m_i)_i^{\mathsf T}\longmapsto\sum_i p_i m_i,
\]
其逆为 \(m\mapsto u_1\otimes(m,0,\ldots,0)^{\mathsf T}\)。
\end{proposition}
\begin{proof}
结合律给出两种矩阵配对的平衡性与外侧作用相容性。为核对第一同构的逆，将行列按标准基展开，得到
\[
(p_i u_i)\otimes(v_jq_j)
=u_1\otimes (p_iE_{1i})(v_jq_j)
=\delta_{ij}(p_iq_j)u_1\otimes v_1.
\]
第一步在 \(S\) 上移动 \(p_iE_{1i}\)，最后一步移动 \((p_iq_j)E_{11}\)。因此全张量等于 \((pq)u_1\otimes v_1\)，两复合均为恒等。第二同构的逆把 \(E_{ij}\) 送到 \(v_i\otimes u_j\)；对任意纯张量，\(R\) 上的平衡关系给出
\[
q\otimes p=\sum_{i,j}v_i(q_i p_j)\otimes u_j,
\]
恰为外积矩阵 \(qp\) 的逆像。

最后，任意行 \(p\) 可写成 \(u_1B\)，其中 \(B\) 的首行为 \(p\)、其余行为零。在张量中把 \(B\) 移到 \(M^n\) 上，得到
\[
p\otimes(m_i)_i^{\mathsf T}
=u_1\otimes\left(\sum_i p_i m_i,0,\ldots,0\right)^{\mathsf T}.
\]
这就验证了所列余单位及其逆在任意 \(M\) 上的两个复合。所有系数保持原次序，证明不要求 \(R\) 交换。
\end{proof}
